%MB9T5
\setcounter{paragraph}{0}
\setcounter{ex}{0}
\PhanI
\loai{Quy luật}
\begin{ex}%[3P2K9-1][Loai1]
Trên một dây có sóng dừng mà các tần số trên dây theo quy luật: $f_1:f_2:f_3:...:f_n=1:2:3:...:n$ trên dây thì
\choice
{số nút bằng số bụng trừ 1}
{\True số nút bằng số bụng cộng 1}
{số nút bằng số bụng}
{số nút bằng số bụng trừ 2}
\loigiai{
\begin{itemize}
\item Nếu sóng dừng trên dây một đầu cố định một đầu tự do thì các tần số $f_1$, $3f_1$, $5f_1$,...
\item Nếu sóng dừng trên dây hai đầu cố định thì các tần số $f_1$, $2f_1$, $3f_1$,...
\item Như vậy, trong bài toán này thì sợi dây hai đầu cố định nên số nút bằng số bụng cộng 1. 
\end{itemize}
}
\end{ex} \haidong{20}
\begin{ex}%[3P2K9-1][Loai1]
Một sợi dây đàn hồi được căng giữa hai điểm cố định. Người ta tạo ra sóng dừng trên dây với tần số bé nhất là $f_1$. Để lại có sóng dừng, phải tăng tần số tối thiểu đến giá trị $f_2$. Tỉ số $\dfrac{f_2}{f_1}$ bằng
\choice
{6}
{4}
{\True 2}
{3}
\loigiai{
Đối với hai đầu cố định thì các tần số tạo ra sóng dừng là 
\begin{align*}
f_{1\min}=\dfrac{v}{2\ell};\ f_2=2\dfrac{v}{2\ell};\ f_3=3\dfrac{v}{2\ell};...\ f_n=n\dfrac{v}{2\ell}
\Rightarrow \dfrac{f_2}{f_1} = 2.
\end{align*}
}
\end{ex} \haidong{20}
\begin{ex}%[3P2K9-1][Loai1]
Một sợi dây đàn hồi được treo thẳng đứng vào một điểm cố định, đầu dưới của dây để tự do. Người ta tạo sóng dừng trên dây với tần số bé nhất là $f_1$. Để có sóng dừng trên dây phải tăng tần số tối thiểu đến giá trị $f_2$. Tỉ số $\dfrac{f_2}{f_1}$ là
\choice
{1,5}
{2}
{2,5}
{\True 3}
\loigiai{
Đối với một đầu cố định một đầu tự do thì các tần số tạo ra sóng dừng là:
\begin{align*}
f_{1\min}=\dfrac{v}{4\ell};\ f_2=3\dfrac{v}{4\ell};\ f_3=5\dfrac{v}{4\ell};... f_n=\left( 2n-1 \right).\dfrac{v}{4\ell } \Rightarrow \dfrac{f_2}{f_1} = 3.
\end{align*}
}\end{ex} \haidong{20}
\loai{Biết hai tần số liên tiếp, tính tần số nhỏ nhất tạo ra sóng dừng}
\begin{ex}%[3P2K9-1][Loai1]
Một sợi dây căng giữa hai điểm cố định. Người ta tạo sóng dừng trên dây. Hai tần số gần nhau nhất cùng tạo ra sóng dừng trên dây là 150 Hz và 200 Hz. Tần số nhỏ nhất tạo ra sóng dừng trên dây đó là
\choice
{\True 50 Hz}
{25 Hz}
{75 Hz}
{100 Hz}
\loigiai{
\begin{itemize}
\item Hiệu hai tần số liên tiếp để dây có sóng dừng là: $\Delta f=\dfrac{v}{2\ell } = 200 – 150 = 50$ Hz.
\item Tần số nhỏ nhất để tạo ra sóng dừng hai đầu cố định là: $f_{\min}=\dfrac{v}{2\ell }=50$ Hz.
\end{itemize}
}
\end{ex} \haidong{20}
\begin{ex}%[3P2K9-1][Loai1]
Một sợi dây đàn hồi có một đầu cố định, một đầu tự do. Thay đổi tần số dao động của sợi dây thì thấy trên dây có sóng dừng với hai tần số liên tiếp là 30 Hz và 50 Hz. Tần số nhỏ nhất để có sóng dừng trên dây là
\choice
{15 Hz}
{20 Hz}
{\True 10 Hz}
{30 Hz}
\loigiai{
\begin{itemize}
\item Ta có Sóng dừng một đầu cố định một đầu thả tự do khi đó chiều dài dây thoả mãn.
\item $\ell = \dfrac{(2k + 1)\lambda}{4} = \dfrac{(2k + 1)v}{4f} \Rightarrow f = \dfrac{(2k + 1)v}{4\ell} \Rightarrow f_{min} = \dfrac{v}{4\ell}; k = 0$.
\item Hai tần số sóng liên tiếp tạo sóng dừng là 30 và 50.
\item $ f_{n + 1}-f_{n} = 20 \Rightarrow \dfrac{2v}{4\ell} = 20 \Rightarrow f_{min} = 10$.
\end{itemize}
}
\end{ex} \haidong{20}
\loai{Chưa biết hai đầu cố định hay tự do}
\begin{ex}%[3P2K9-1][Loai1]
Một sợi dây đàn hồi dài có sóng dừng với hai tần số liên tiếp là 30 Hz và 50 Hz. Phát biểu nào sau đây là đúng?
\choice
{Dây đó một đầu cố định và một đầu tự do. Tần số nhỏ nhất để có sóng dừng khi đó là 30 Hz}
{\True Dây đó một đầu cố định và một đầu tự do. Tần số nhỏ nhất để có sóng dừng khi đó là 10 Hz}
{Dây đó có hai đầu cố định. Tần số nhỏ nhất để có sóng dừng khi đó là 30 Hz}
{Dây đó có hai đầu cố định. Tần số nhỏ nhất để có sóng dừng khi đó là 10 Hz}
\loigiai{
\begin{itemize}
\item Xét tỉ số $\dfrac{f_1}{f_2}=\dfrac{30}{50}=\dfrac{3}{5}$ là hai số lẻ liên tiếp nên sợi dây có một đầu cố định một đầu tự do.
\item Tần số nhỏ nhất để tạo sóng dừng: $f_{\min} =0,5\left|f_1-f_2\right|=10\ Hz$. 
\end{itemize}
}
\end{ex} \haidong{20}
\loai{Hai tần số liên tiếp cùng tạo ra sóng dừng. Tìm tốc độ truyền sóng}
\begin{ex}%[3P2K9-1][Loai1]
Một sợi dây căng giữa hai điểm cố định cách nhau 60 cm. Hai sóng có tần số gần nhau liên tiếp cùng tạo ra sóng dừng trên dây là 84 Hz và 98 Hz. Biết tốc độ truyền của các sóng trên dây là bằng nhau. Tốc độ truyền sóng trên dây là
\choice
{1,44 m/s}
{1,68 m/s}
{\True 16,8 m/s}
{14,4 m/s}
\loigiai{
\begin{itemize}
\item Ta có: $\ell = n\dfrac{\lambda}{2} = k\dfrac{v}{2f} \Rightarrow f = k\dfrac{v}{2\ell}= kf_0$, với $f_0 = \dfrac{v}{2\ell}$.
\item $ \Delta f = f_0 = 98-84 = 14\ Hz\Rightarrow v = 2\ell f_0 = 2.0,6.14 = 16,8\ m/s$.
\end{itemize}
}
\end{ex} \haidong{20}

\loai{Điều kiện kẹp tần số- Tìm tần số}
\begin{ex}%[3P2K9-1][Loai1]
Cho một thanh thép mảnh, thẳng, dài 3 m, một đầu bị kẹp chặt, đầu còn lại để tự do. Kích thích thanh thép dao động theo phương vuông góc làm hình thành sóng dừng. Tốc độ truyền sóng trên thanh thép bằng 1100 m/s. Biết tần số của sóng trong khoảng từ 600 Hz đến 734 Hz. Tần số sóng truyền trên thanh khi đó xấp xỉ bằng
\choice
{720 Hz}
{\True 642 Hz}
{713Hz}
{651Hz}
\loigiai{
Ta có: $ \ell = \left(k + \dfrac{1}{2}\right).\dfrac{\lambda}{2} = \dfrac{2n + 1}{4}.\dfrac{v}{f}\Rightarrow f = \left(2k + 1\right) \dfrac{v}{4\ell}$\\
$\Rightarrow 600 < \left(2k + 1\right) \dfrac{1100}{4.3} < 734\Rightarrow 2,77 < k < 3,5\Rightarrow k = 3\Rightarrow f = (2.3 + 1)\dfrac{1100}{4.3} = 642\ Hz$.
}
\end{ex} \haidong{20}
\loai{Điều kiện kẹp tần số- Tìm số lần tạo sóng dừng}
\begin{ex}%[3P2K9-1][Loai1]
Một sợi dây đàn hồi dài 1 m được treo lơ lửng lên một cần rung. Cần có thể rung theo phương ngang với tần số thay đổi được từ 100 Hz đến 120 Hz. Vận tốc truyền sóng trên dây 8 m/s. Mỗi khi hình thành sóng dừng thì đầu dây gắn với cần rung được coi là một nút sóng. Hỏi trong quá trình thay đổi tần số rung của cần, số lần có sóng dừng trên dây là
\choice
{10 lần}
{4 lần}
{\True 5 lần}
{12 lần}
\loigiai{
\begin{itemize}
\item Ta có dây đàn hồi được treo lơ lửng trên một cần rung như vậy ứng với sóng dừng một đầu cố định một đầu tự do.
\item Khi đó $\ell = \dfrac{(2k + 1)\lambda}{4} = \dfrac{(2k + 1)v}{4f} \Rightarrow f = \dfrac{(2k + 1)v}{4\ell} = 2(2k + 1)$.
\item Theo đề: $100 \leq f \leq 120 \Rightarrow 24,5 \leq k \leq 29,5 \Rightarrow$ Có 5 giá trị k thoả mãn.
\item Như vậy có thể tạo được 5 lần sóng dừng trên dây khi thay đổi tần số rung của cần.
\end{itemize}
}
\end{ex} \haidong{20}

\begin{ex}%[3P2K9-1][Loai1]
Một sợi dây đàn hồi một đầu cố định, một đầu gắn với âm thoa có tần số thay đổi được. Khi thay đổi tần số âm thoa thấy với hai giá trị liên tiếp của tần số là 28 Hz và 42 Hz thì trên dây có sóng dừng. Coi vận tốc truyền sóng và chiều dài dây là không đổi. Nếu tăng dần giá trị tần số từ 0 Hz đến 50 Hz sẽ có bao nhiêu giá trị của tần số để trên dây lại có sóng dừng?
\choice
{7 giá trị}
{6 giá trị}
{4 giá trị}
{\True 3 giá trị}
\loigiai{
\begin{itemize}
\item Vì sợi dây hai đầu cố định nên
$f_{\min}=f_{k+1}-f_k=42-28=14\ Hz\Rightarrow f_k=14k$.
\item Thay vào điều kiện $0<f<50\ Hz\Rightarrow 0<k\le 3,5\Rightarrow k=1;\ 2;\ 3 \Rightarrow $ có ba lần tạo sóng dừng.
\end{itemize}
}
\end{ex} \haidong{20}
\loai{Số nút, bụng thay đổi}
\begin{ex}%[3P2K9-1][Loai1]
Khi có sóng dừng trên sợi dây đàn hồi AB thì thấy trên dây có 7 nút (kể cả 2 nút ở hai đầu A, B) với tần số sóng là 42 Hz. Cũng với dây AB và tốc độ truyền sóng như trên, muốn trên dây có 5 nút (tính cả hai đầu A, B) thì tần số sóng có giá trị là
\choice
{30 Hz}
{63 Hz}
{\True 28 Hz}
{58,8 Hz}
\loigiai{
\begin{itemize}
\item 7 nút $\Rightarrow$  số bụng là 6: $42=6\dfrac{v}{2\ell }$.
\item 5 nút $\Rightarrow$  số bụng là 4: $f=4\dfrac{v}{2\ell } \Rightarrow  f = 28\ Hz$.
\end{itemize}}
\end{ex} \haidong{20}
\begin{ex}%[3P2K9-1][Loai1]
Sóng truyền trên một sợi dây có một đầu cố định, một đầu tự do. Khi tần số sóng trên dây là 190 Hz, trên dây hình thành sóng dừng với 10 bụng sóng. Tốc độ truyền sóng không thay đổi. Tần số nào cho dưới đây cũng tạo ra sóng dừng trên dây?
\choice
{20 Hz}
{40 Hz}
{\True 50 Hz}
{100 Hz}
\loigiai{
\begin{itemize}
\item Sóng dừng một đầu cố định một đầu tự do: 190 = $\left( 10. 2-1 \right). \dfrac{v}{4\ell }=19{{f}_{\min }} \Rightarrow f_{\min} = 10$ Hz.
\item Điều kiện tần số để xảy ra sóng dừng là $f=\left( 2n-1 \right)\dfrac{v}{4\ell }=\left( 2n-1 \right). 10$: lẻ lần của 10.
\end{itemize}
}
\end{ex} \haidong{20}

\begin{ex}%[3P2K9-1][Loai1]
Một thanh thép mảnh, thẳng, dài, treo ngang trên các sợi chỉ nhẹ, hai đầu để tự do. Kích thích để trên thanh thép tạo thành sóng dừng (sóng ngang) với tần số $f=2560$ Hz thì thấy trên thanh thép có hai điểm không dao động. Để hình thành sóng dừng trên thanh thép với chỉ một điểm không dao động thì phải kích thích thanh thép dao động với tần số bằng
\choice
{1180 Hz}
{853 Hz}
{\True 1280 Hz}
{1507 Hz}
\loigiai{
\begin{itemize}
\item Khi trên thanh thép có hai điểm không dao động thì $\ell=2.\dfrac{\lambda}{2} = \dfrac{v}{f}\quad (1)$
\item Khi trên thanh thép có 1 điểm không dao động thì $\ell=\dfrac{\lambda}{2} = \dfrac{1}{2}\dfrac{v}{f'}\quad (2)$
\item Từ $(1)$ và $(2) \Rightarrow f' = \dfrac{f}{2} = \dfrac{2560}{2} = 1280\ Hz$.
\end{itemize}
}
\end{ex} \haidong{20}

\loai{Số nút, bụng thay đổi kèm tần số}
\begin{ex}%[3P2K9-1][Loai1]
Một sợi dây AB dài 4,5 m có đầu dưới A để tự do, đầu trên B gắn với một cần rung với tần số f có thể thay đổi được. Ban đầu trên dây có sóng dừng với đầu A bụng đầu B nút. Khi tần số f tăng thêm 3 Hz thì số nút trên dây tăng thêm 18 nút và A vẫn là bụng B vẫn là nút. Tính tốc độ truyền sóng trên sợi dây.
\choice
{3,2 m/s}
{1,0 m/s}
{\True 1,5 m/s}
{3,0 m/s}
\loigiai{
\begin{itemize}
\item Ta có: $\Delta f=\Delta k\dfrac{v}{2\ell}\Rightarrow 3=18.\dfrac{v}{2.4,5}\Rightarrow v=1,5\ m/s$.
\end{itemize}
}
\end{ex} \haidong{20}
\begin{ex}%[3P2K9-1][Loai1]
Một sợi dây CD dài 1 m, đầu C cố định, đầu D gắn với cần rung với tần số thay đổi được. D được coi là nút sóng. Ban đầu trên dây có sóng dừng. Khi tần số tăng thêm 20 Hz thì số nút trên dây tăng thêm 7 nút. Sau khoảng thời gian bằng bao nhiêu sóng phản xạ từ C truyền hết một lần chiều dài sợi dây?
\choice
{\True 0,175 s}
{0,07 s}
{0,5 s}
{1,2 s}
\loigiai{
\begin{itemize}
\item Gọi số bụng ban đầu là k , số bụng sóng lúc sau là $k+7$.
\item Ta có: $\ell = k\dfrac{\lambda}{2} = \dfrac{kv}{2f} = \dfrac{(k + 7)v}{(f + 20)} \Rightarrow \dfrac{7}{k} = \dfrac{20}{f} \Rightarrow k = \dfrac{7f}{20}$.
\item Từ đó ta có $\dfrac{7v}{40} = \ell\Rightarrow$ thời gian sóng truyền hết một lần sợi dây: $t = \dfrac{\ell}{v} = \dfrac{7}{40} = 0,175$ s.
\end{itemize}
}
\end{ex} \haidong{20}

\loai{Thay đổi tần số lượng nhỏ nhất, không thay đổi tính chất 2 đầu}

\begin{ex}%[3P2K9-1][Loai1]
Trên một sợi dây đàn hồi dài 45 cm buông thẳng đứng, đầu trên cố định, đầu dưới tự do đang có một sóng dừng với tần số xác định. Để trên dây tiếp tục hình thành sóng dừng thì phải thay đổi tần số kích thích một lượng nhỏ nhất bằng 50 Hz. Tốc độ truyền sóng trên dây bằng
\choice
{75 m/s}
{\True 45 m/s}
{7,5 m/s}
{15 m/s}
\loigiai{
\begin{itemize}
\item Ta có: $\ell = \left(k + \dfrac{1}{2}\right)\dfrac{v}{2\ell} = (2k + 1)\dfrac{v}{4f}$.
\item $ f = (2k + 1)\dfrac{v}{4\ell} = (2k + 1)f_0$ với $f_0 = \dfrac{v}{4\ell}\Rightarrow \Delta f = 2f_0 = 50 \Rightarrow f_0 = 25\ Hz$.
\item $ v = 4\ell.f_0 = 4.0,45.25 = 45\ m/s$.
\end{itemize}
}
\end{ex} \haidong{20}
\begin{ex}%[3P2K9-1][Loai1]
Trên một sợi dây đàn hồi dài 55 cm buông thẳng đứng, đầu trên cố định, đầu dưới tự do đang có một sóng dừng với tần số xác định. Để trên dây tiếp tục hình thành sóng dừng thì phải thay đổi tần số kích thích một lượng nhỏ nhất bằng 30 Hz. Tốc độ truyền sóng trên dây bằng
\choice
{25 m/s}
{\True 33 m/s}
{43 m/s}
{15 m/s}
\loigiai{
\begin{itemize}
\item Ta có: $\ell = \left(n + \dfrac{1}{2}\right)\dfrac{v}{2\ell} = \left(2n + 1\right)\dfrac{v}{4f}\Rightarrow f = \left(2n + 1\right)\dfrac{v}{4\ell} = \left(2n + 1\right)f_0$ với $f_0 = \dfrac{v}{4\ell}$.
\item $ \Delta f = 2f_0 = 30 \Rightarrow f_0 = 15\ Hz\Rightarrow v = 4\ell.f_0 = 4.0,55.15 = 33\ m/s$.
\end{itemize}
}\end{ex} \haidong{20}

\loai{Thay đổi đầu cố định tự do. Tần số thay đổi lượng nhỏ nhất}
\begin{ex}%[3P2G9-1][Loai1]
Một sợi dây đàn hồi, đầu A gắn với nguồn dao động và đầu B tự do. Khi dây rung với tần số $f = 12$ Hz thì trên dây xuất hiện sóng dừng ổn định có 8 điểm nút trên dây với A là nút và B là bụng. Nếu đầu B được giữ cố định và tốc độ truyền sóng trên dây không đổi thì phải thay đổi tần số rung của dây một lượng nhỏ nhất bằng bao nhiêu để trên dây tiếp tục xảy ra hiện tượng sóng dừng ổn định?
\choice
{$\dfrac{4}{3}$ Hz}
{\True 0,8 Hz}
{12 Hz}
{1,6 Hz}
\loigiai{
\begin{itemize}
\item Lúc đầu, một đầu cố định một đầu tự do thì trên dây có sóng dừng với tần số f:\\
$\ell=\left(2n+1\right)\dfrac{\lambda}{4}=\left(2n+1\right)\dfrac{v}{4f}\Rightarrow \dfrac{v}{2\ell}=\dfrac{2f}{\left(2n+1\right)}$. 
\item Vì số nút $=$ số bụng $= n +1= 8$ nên:
$\dfrac{v}{2\ell}=\dfrac{2.12}{\left(2.7+1\right)}=1,6\ Hz$.
\item Sau đó, giữ đầu cố định hai đầu thì trên dây có sóng dừng với tần số f':\\
$\ell=k\dfrac{\lambda}{2}=k\dfrac{v}{2f'}\Rightarrow f'=k\dfrac{v}{2\ell}=1,6k$.
\item Độ biến thiên tần số: $\Delta f=\left|{f'-f}\right|=\left|{1,6k-12}\right|\ Hz$.
\item Để tìm $\Delta f_{\min}$ ta cho $\Delta f=0\Rightarrow k=7,5.$ 
\item Nhưng vì k nguyên nên $k = 7$ hoặc $k = 8$.
\item Do đó, $\Delta f_{\min}=\left|{1,6.7-12}\right|=0,8\ Hz$.
\item \textbf{Chú ý:} Đến đây ta rút ra công thức giải nhanh: $\Delta f_{\min}=\dfrac{f}{\left(2n+1\right)}=\dfrac{f'_{\min}}{2}.$
\end{itemize}
}
\end{ex} \haidong{20}




\setcounter{paragraph}{0}
\setcounter{ex}{0}
\PhanII
\begin{ex}
Một sợi dây đàn hồi dài $90 ~cm$ đang có sóng dừng với hai tần số liên tiếp là $30 ~Hz$ và $50 ~Hz$. Biết tốc độ truyền sóng là $40 ~cm/s$.
\choiceTF[t]
{\True Dây có một đầu cố định và một đầu tự do}
{Tần số nhỏ nhất để có sóng dừng khi đó là $30 ~Hz$}
{Bước sóng của sợi dây khi đang có sóng dừng với tần số nhỏ nhất là $15 ~cm$}
{Với tần số nhỏ nhất để có sóng dừng trên dây thì trên dây có 10 nút sóng}
\loigiai{
\begin{itemchoice}
\itemch 
Gọi $f_{n}=30 ~Hz$ và $f_{n+1}=50 ~Hz$.\\
$\dfrac{f_{n+1}}{f_{n}}=\dfrac{5}{3}$\\
Vậy sóng dừng xảy ra trên dây ứng với trường hợp một đầu cố định, một đầu tự do.
\itemch Tần số nhỏ nhất để có sóng dừng là $f_{\text {min}}=\dfrac{|30-50|}{2}=10 ~Hz$
\itemch Vì $\lambda=\dfrac{v}{f_{\text {min}}}=\dfrac{40}{10}=4 ~cm$
\itemch Vì $k=\dfrac{f_{n}}{f_{\text {min}}}=\dfrac{30}{10}=3$
\end{itemchoice}
}
\end{ex} \haidong{20}
\begin{ex}
Một dây đàn guitar dài $L = 0,64\ m$ (tính từ lược đàn đến ngựa đàn) tạo sóng dừng cơ bản khi được gảy, phát ra âm có tần số $f_{1}=83,5\ Hz$. Biết tốc độ truyền âm trong không khí $v_{kk}=343,2\ m/s$.
\immini{\choiceTF[t]
{\True Bước sóng của âm trong không khí là $\lambda_{kk}=4,11\ m$}
{\True Vận tốc truyền âm trên dây đàn là $v_d = 106,9\ m/s$}
{\True Âm truyền trong không khí nhanh hơn trong dây đàn xấp xỉ $3,2$ lần}
{Ở chế độ sóng dừng cơ bản, chiều dài dây đàn đúng bằng một bước sóng của sóng trên dây}}{
\includegraphics[width=5 cm]{images/2025_04_24_d0a7e0ec0735763b7a55g-5}}
\loigiai{
\begin{itemchoice}
\itemch %$\lambda_{kk}=\dfrac{v_{kk}}{f_{1}}=\dfrac{343,2}{83,5}\approx 4,11\ m$
\itemch Ở chế độ cơ bản: $\lambda_d=2L=2\  0,64=1,28. m$, $v_d=\lambda_d f_{1}=1,28. 83,5\approx 106,9\ m/s$
\itemch $\dfrac{v_{kk}}{v_d}=\dfrac{343,2}{106,9}\approx3,2$
\itemch Chế độ cơ bản có $L=\dfrac{\lambda_d}{2}$, không phải $L=\lambda_d$
\end{itemchoice}
}
\end{ex}
\begin{ex}
Hình vẽ là sơ đồ thí nghiệm tạo sóng dừng trên dây. Sợi dây căng ngang có chiều dài $L$, một đầu cố định, một đầu gắn với cần rung. Tần số cần rung được điều chỉnh từ $10\ Hz$ đến $100\ Hz$ bởi máy phát tần số. Các bước tiến hành thí nghiệm như sau:
\begin{itemize}
\item 	Bước 1: Bật máy phát tần số.	
\item Bước 2: Tăng giá trị tần số. Khi trên dây xuất hiện ổn định những vị trí đứng yên và những vị trí dao động với biên độ lớn nhất thì ghi lại số bụng sóng $n$ và giá trị tần số $f$ vào bảng.
\end{itemize}
\begin{center}
\begin{minipage}{0.48\textwidth}
\centering
\includegraphics[width=6cm]{images/2025_04_24_93aaf10d3a9ed4c40f34g-4}
\end{minipage}\hfill
\begin{minipage}{0.48\textwidth}
\centering
\renewcommand{\arraystretch}{1.2}% giãn dòng
\begin{tabular}{|c|c|}
\hline
\textbf{Số bụng sóng $n$} & \textbf{Tần số $f\ (Hz)$}\\
\hline
2 & 26\\
\hline
3 & 39\\
\hline
4 & 52\\
\hline
$x$ & 78\\
\hline
\end{tabular}
\end{minipage}
\end{center}

\choiceTF[t]
{Khi tần số máy phát là $12\ Hz$ thì trên dây có sóng dừng với $1$ bụng sóng}
{Giá trị của $x$ là $5$}
{Công thức liên hệ giữa tần số khi có sóng dừng và số bụng sóng là $f=0,5nL$ với n là số nguyên}
{\True Nếu tiếp tục tăng tần số từ 79 Hz đến 150 Hz, sẽ có 5 tần số cho sóng dừng trên dây}
\loigiai{
\begin{itemchoice}
\itemch Sai
\itemch Sai
\itemch Sai
\itemch Đúng
\end{itemchoice}
}
\end{ex}
\begin{ex}
Hai nguồn dao động cùng biên độ, cùng pha, cùng tần số $f=10\ Hz$, đặt tại hai điểm A, B trên mặt nước tạo ra hai sóng lan truyền với tốc độ $v=60\ cm/s$.
\choiceTF[t]
{\True Phần tử tại M trên mặt nước cách A, B những khoảng $MA=24\ cm$ và $MB=36\ cm$ luôn dao động với biên độ cực đại}
{Trên đoạn thẳng AB, hai phần tử gần nhau nhất không dao động cách nhau 6 cm}
{\True Phần tử mặt nước tại trung điểm O của AB dao động với biên độ cực đại}
{Dùng nguồn dao động nói trên nối với một đầu của một sợi dây đàn hồi, đầu còn lại của dây được giữ cố định thì trên dây xuất hiện sóng dừng với hai bụng sóng. Để trên dây có sóng dừng với 5 bụng sóng thì phải tăng tần số của nguồn thêm 25 Hz}
\loigiai{
\begin{itemchoice}
\itemch Đúng.
\itemch Sai.
\itemch Đúng.
\itemch Sai.
\end{itemchoice}
}
\end{ex}















\setcounter{paragraph}{0}
\setcounter{ex}{0}
\PhanIII
\begin{ex}
Trong hiện tượng sóng dừng, nguồn dao động có tần số thay đổi được gây ra sóng lan truyền trên dây một đầu cố định, một đầu tự do. Thay đổi tần số của nguồn thì nhận thấy có hai tần số liên tiếp $30 ~Hz$ và $50 ~Hz$ trên dây hình thành sóng dừng. Để sóng hình thành trên đây với 4 bụng sóng thì tần số của nguồn dao động là bao nhiêu Hz?\\\shortans[oly]{70}
\loigiai{
70
}
\end{ex} \haidong{20}
\begin{ex}
Một sợi dây căng giữa hai điểm cố định cách nhau $80 ~cm$. Hai sóng có tần số gần nhau liên tiếp cùng tạo ra sóng dừng trên dây là $70 ~Hz$ và $84 ~Hz$. Biết tốc độ truyền sóng trên dây không đổi. Tìm tốc độ truyền sóng trên dây (đơn vị m/s).\\\shortans[oly]{22,4}
\loigiai{
22,4
}
\end{ex} \haidong{20}


\loai{Thay đổi đầu cố định tự do. Tìm tần số, tốc độ truyền sóng}
\begin{ex}%[3P2K9-1][Loai1]
Quan sát sóng dừng trên sợi dây AB, đầu A dao động điều hòa theo phương vuông góc với sợi dây (coi A là nút). Với đầu B tự do và tần số dao động của đầu A là 22 Hz thì trên dây có 6 nút. Nếu đầu B cố định và coi tốc độ truyền sóng trên dây như cũ, để vẫn có 6 nút thì tần số dao động của đầu A phải bằng bao nhiêu héc?
\shortans[oly]{20}
\loigiai{
\begin{itemize}
\item B tự do: 6 nút $\Rightarrow$ số bụng cũng là 6: $22=(2. 5+1)\dfrac{v}{4\ell }$.
\item B cố định: 6 nút $\Rightarrow$ 5 bụng: $f=5. \dfrac{v}{2\ell }\Rightarrow f = 20$ Hz.
\end{itemize}
}
\end{ex} \haidong{20}
\begin{ex}
\immini[thm]{Trong thí nghiệm đo tốc độ truyền âm trong không khí bằng hiện tượng sóng dừng với nguồn âm có tần số $500 \  Hz$ như hình vẽ bên. Khi di chuyển nút cao su bên trong ống thủy tinh người ta thấy tại ba vị trí liên tiếp thước đo có giá trị là $34 \  cm, 68 \  cm$ và $102 \  cm$ thì âm phát ra lớn nhất. Tốc độ truyền âm đo được trong thí nghiệm có giá trị là bao nhiêu $m/s$?} 	{\includegraphics[width=8 cm]{images/2025_04_24_6110796b3ccfb0c2ffbfg-3}}
\shortans[oly]{340}
\loigiai{

}
\end{ex}
\loai{$f_{\min}$ - Thay đổi $\ell$}
\begin{ex}%[3P2K9-1][Loai1]
Một sợi dây đàn hồi một đầu cố định, một đầu tự do. Tần số dao động bé nhất để sợi dây có sóng dừng là $f_0$. Tăng chiều dài thêm 1 m thì tần số dao động bé nhất để sợi dây có sóng dừng là 5 Hz. Giảm chiều dài bớt 1 m thì tần số dao động bé nhất để sợi dây có sóng dừng là 20 Hz. Giá trị của $f_0$ là bao nhiêu héc?
\shortans[oly]{8}
\loigiai{
\begin{itemize}
\item Vì sợi dây một đầu cố định và một đầu tự do nên điều kiện sóng dừng là
\begin{align*}\ell=\left(2k-1\right)\dfrac{\lambda}{4}=\left(2k-1\right)\dfrac{v}{4\ell}\Rightarrow f_k=\left(2k-1\right)\dfrac{v}{4\ell}\Rightarrow f_{\min}=\dfrac{v}{4\ell}. \end{align*}
\item Áp dụng công thức này cho hai trường hợp:
$\begin{cases}
5=\dfrac{v}{4(\ell+1)}\\
20=\dfrac{v}{4\left(\ell-1\right)}\\
\end{cases}\Rightarrow \begin{cases}
\ell=\dfrac{5}{3}\ m.\\
v=\dfrac{160}{3}\ m/s.\\
\end{cases}$
\item $f_0=f_{\min}=\dfrac{v}{4\ell}=\dfrac{\dfrac{160}{3}}{4.\dfrac{5}{3}}=8\ Hz$.
\end{itemize}
}
\end{ex} \haidong{20}











\begin{ex}%[3P2T9-1][Loai1]
Một dây đàn hồi căng ngang, hai đầu cố định. Thấy hai tần số tạo ra sóng dừng trên dây là 2964 Hz và 4940 Hz. Biết tần số nhỏ nhất tạo ra sóng dừng nằm trong khoảng từ 380 Hz đến 720 Hz. Với tần số nằm trong khoảng từ 8 kHz đến 11 kHz, có bao nhiêu tần số tạo ra sóng dừng?
\shortans[oly]{6}
\loigiai{
\begin{itemize}
\item Luôn có: $f=n\dfrac{v}{2\ell }=n. {{f}_{\min }}$. Điều kiện: $380\ Hz < f_{\min} < 720\ Hz\quad (*)$
\item Ở tần số 2964 Hz, giả sử có m bụng sóng $\Rightarrow 2964 = m$. $f_{\min}\ (1) \Rightarrow f_{\min} = \dfrac{2964}{m}$. 
\item Từ $(*) \Rightarrow 4,12 < m < 7,8$.
\item Ở tần số 4940 Hz, giả sử có k bụng sóng $\Rightarrow 4940 = k;\ f_{\min}\ (2)\Rightarrow f_{\min} = \dfrac{4940}{k}$. 
\item Từ $(*) \Rightarrow 6,86 < k < 13$.
\item Chia từng vế $(1)$ với $(2) \Rightarrow \dfrac{m}{k}=\dfrac{3}{5}=\dfrac{6}{10}=\dfrac{9}{15}=... $
\item So sánh với điều kiện của m và k $\Rightarrow (m,k) = (6,10)$.
\item $f_{\min} = 494$ Hz $\Rightarrow$ Tần số để xảy ra sóng dừng phải là: $f = 494n$ Hz.
\item Tần số xảy ra sóng dừng trong khoảng từ 8 kHz đến 11 kHz là: \\$8000 < 494n < 11000 \Rightarrow 16,2 < n < 22,3$.
\item $n = 17, 18, 19, 20, 21, 22$: có 6 tần số trong khoảng đó có thể tạo ra sóng dừng.
\end{itemize}
}
\end{ex} \haidong{20}
